%! Characteristics of Linear Functions — Unit 1, Lesson 1.0 — Group Activity (Answer Key)
% Mirrors activity/main.tex; each \writelines{2} is answered with exactly two \ansline's so the
% blank and the key occupy the same height.
\documentclass[10pt]{article}
\usepackage{algebra2-article}
\usepackage{algebra2-key}   % pulls in -boxes; do NOT also load -boxes
% Temperature graph: hours h on the horizontal axis (0..7), temperature T on the vertical (-8..8).
\newcommand{\tempgrid}{%
  \draw[step=1, linegray!40, very thin] (0,-8) grid (7,8);
  \draw[->, semithick, charcoal] (0,0)--(7.6,0) node[right, font=\scriptsize]{$h$};
  \draw[->, semithick, charcoal] (0,-8.6)--(0,8.8) node[above, font=\scriptsize]{$T$ ($^\circ$C)};
  \foreach \x in {1,...,7} \node[charcoal, font=\tiny, below] at (\x,0){\x};
  \foreach \y in {-8,-6,-4,-2,2,4,6,8} \node[charcoal, font=\tiny, left] at (0,\y){\y};}

\begin{document}

\pageheader{Unit 1, Lesson 1.0}{Group Activity --- Answer Key}

\vspace{0.08in}

\begin{headlinebox}{forestbg}
\small\textbf{Freezing Point.} On a cold weekend the number that matters most is \emph{zero}: above it
the roads are wet, below it they are ice. Work with your group. Use the names from the notes ---
\textbf{slope}, \textbf{$y$-intercept}, \textbf{zero}, \textbf{increasing/decreasing},
\textbf{positive/negative} --- and for every answer be ready to show \emph{where you see it on the
graph}.
\end{headlinebox}

\vspace{0.06in}

\begin{scenariobox}[1. Saturday Morning]{forest}
At 6:00 a.m.\ on Saturday the temperature in Blacksburg was $-6\,^\circ$C. It rose steadily,
$2\,^\circ$C every hour, all morning (until noon). Let $h$ be the number of hours after 6:00 a.m.
and $T$ the temperature in $^\circ$C.

\begin{enumerate}[label=\textbf{\alph*.}, leftmargin=1.8em, itemsep=5pt]
  \item Complete the table.
    \begin{center}
    \renewcommand{\arraystretch}{1.4}
    \begin{tabular}{|c|c|c|c|c|c|c|c|}
    \hline
    $h$ & $0$ & $1$ & $2$ & $3$ & $4$ & $5$ & $6$ \\ \hline
    $T$ & $-6$ & $-4$ & \ans{$-2$} & \ans{$0$} & \ans{$2$} & \ans{$4$} & \ans{$6$} \\ \hline
    \end{tabular}
    \end{center}
  \item Write a rule for the temperature $h$ hours after 6:00 a.m. \quad
        $T(h) = $ \ans{$2h-6$}
\end{enumerate}

\begin{minipage}[t]{0.58\linewidth}
\vspace{0pt}
\begin{enumerate}[label=\textbf{\alph*.}, start=3, leftmargin=1.8em, itemsep=5pt]
  \item Circle the \textbf{$y$-intercept} on the graph: $(\ans{$0$},\ \ans{$-6$})$.
        In the story it tells you:
        \ansline{the starting temperature --- at 6:00 a.m.\ it was $-6\,^\circ$C,}
        \ansline{six degrees below freezing.}
  \item Circle the \textbf{zero}: $(\ans{$3$},\ \ans{$0$})$. What time of day is that,
        and why is that moment special?
        \ansline{9:00 a.m. --- the temperature reaches $0\,^\circ$C, freezing.}
        \ansline{Before it the roads are icy; after it, only wet.}
\end{enumerate}
\end{minipage}\hfill
\begin{minipage}[t]{0.38\linewidth}
\vspace{0pt}\centering
\begin{tikzpicture}[x=0.55cm, y=0.24cm]
  \tempgrid
  \draw[very thick, forest] (0,-6)--(7,8);
  \fill[charcoal] (0,-6) circle (2.8pt);
  \fill[charcoal] (3,0) circle (2.8pt);
\end{tikzpicture}
\end{minipage}

\begin{enumerate}[label=\textbf{\alph*.}, start=5, leftmargin=1.8em, itemsep=5pt]
  \item The \textbf{slope} is \ans{$2$} $^\circ$C per hour. Circle where that number sits in
        your rule, then say how you \emph{see} it on the graph from one hour to the next.
        \ansline{It is the coefficient of $h$ in $T(h)=2h-6$.}
        \ansline{On the graph: one step right along $h$ is two units up along $T$.}
  \item Is $T$ \textbf{increasing} or \textbf{decreasing}? \ans{increasing} \;
        How does the \emph{graph} show it (not the story)? \ans{it rises left to right} \\[3pt]
        $T$ is \textbf{negative} when \ans{$0\le h<3$} \qquad
        $T$ is \textbf{positive} when \ans{$h>3$}
\end{enumerate}
\end{scenariobox}

\boxguard[22]
\begin{scenariobox}[2. Sunday Evening]{navy}
At 4:00 p.m.\ on Sunday it was $6\,^\circ$C, and the temperature \emph{fell} steadily, $2\,^\circ$C
every hour, through the evening. Now $h$ counts the hours after 4:00 p.m.

\begin{minipage}[t]{0.56\linewidth}
\vspace{0pt}
\begin{enumerate}[label=\textbf{\alph*.}, leftmargin=1.8em, itemsep=5pt]
  \item Write a rule: \; $T(h) = $ \ans{$6-2h$} \\[2pt]
        Slope: \ans{$-2$} \quad Increasing or decreasing? \ans{decreasing}
  \item The $y$-intercept is $(\ans{$0$},\ \ans{$6$})$ and the zero is
        $(\ans{$3$},\ \ans{$0$})$. What does each one tell you about Sunday evening?
        \ansline{$(0,6)$: at 4:00 p.m.\ it was $6\,^\circ$C --- the start.}
        \ansline{$(3,0)$: at 7:00 p.m.\ it hit freezing --- then ice.}
\end{enumerate}
\end{minipage}\hfill
\begin{minipage}[t]{0.40\linewidth}
\vspace{0pt}\centering
\begin{tikzpicture}[x=0.55cm, y=0.24cm]
  \tempgrid
  \draw[very thick, navy] (0,6)--(7,-8);
  \fill[charcoal] (0,6) circle (2.8pt);
  \fill[charcoal] (3,0) circle (2.8pt);
\end{tikzpicture}
\end{minipage}

\begin{enumerate}[label=\textbf{\alph*.}, start=3, leftmargin=1.8em, itemsep=5pt]
  \item At 5:00 p.m.\ ($h=1$) the temperature is \ans{$4$} $^\circ$C. Is that number
        \textbf{positive} or \textbf{negative}? \ans{positive} \; At that moment, is the
        temperature \textbf{rising} or \textbf{falling}? \ans{falling} \\[3pt]
        So: can a temperature be \emph{positive} and \emph{falling} at the same time? Explain what
        each word is describing.
        \ansline{Yes. \emph{Positive} describes where the graph \emph{sits} --- above the $h$-axis.}
        \ansline{\emph{Falling} describes which way it \emph{moves} --- the slope is negative.}
  \item $T$ is \textbf{positive} when \ans{$0\le h<3$} \qquad
        $T$ is \textbf{negative} when \ans{$h>3$}
  \item \textbf{The story, then the math.} In the story, which hours actually make sense as inputs?
        \ans{the evening hours, $0\le h\le 8$} \\[3pt]
        Now \emph{forget the story}: the rule $T(h)=6-2h$ accepts any number $h$, so its
        \textbf{domain} is \ans{all reals} and its \textbf{range} is \ans{all reals}.
        Why are those two answers different?
        \ansline{The rule has no limits --- its line runs forever both ways.}
        \ansline{The story does: no ``$-3$ hours after 4:00 p.m.,'' and the evening ends.}
\end{enumerate}
\end{scenariobox}

\boxguard[16]
\begin{scenariobox}[3. Side by Side]{forest}
Look at both graphs together. They have the \emph{same} zero at $h=3$.
\begin{enumerate}[label=\textbf{\alph*.}, leftmargin=1.8em, itemsep=5pt]
  \item Saturday is negative \ans{before} $h=3$ and positive \ans{after} it.
        Sunday is the other way around.
  \item Their slopes are \ans{$2$} and \ans{$-2$}. What do the two \emph{signs} tell
        you, and what does the fact that both are $2$ in size tell you?
        \ansline{The signs give direction: Saturday rises, Sunday falls.}
        \ansline{Both are $2$ in size, so both change $2\,^\circ$ per hour --- equally steep.}
  \item Write one sentence, using the words \emph{zero} and \emph{sign}, that explains why neither
        line can switch from negative to positive anywhere except at $h=3$.
        \ansline{A line changes sign only where it crosses the $h$-axis --- its zero.}
        \ansline{Each of these lines has exactly one zero, and it is at $h=3$.}
\end{enumerate}
\end{scenariobox}

\boxguard[22]
\begin{scenariobox}[4. Model It --- The Car-Wash Card]{navy}
A car-wash gift card starts with \$40, and each wash costs \$8, so the balance after $w$ washes is
$B(w) = 40 - 8w$.
\begin{enumerate}[label=\textbf{\alph*.}, leftmargin=1.8em, itemsep=5pt]
  \item What does the slope $-8$ tell you about the card? And the $y$-intercept $(0,40)$?
        \ansline{The slope: every wash takes \$8 off the balance (it falls, so it is negative).}
        \ansline{The $y$-intercept: the card holds \$40 before any wash.}
  \item Find the zero of $B$ by solving $B(w)=0$.
    \begin{work}
      40 - 8w &= 0 \\
          -8w &= -40 \\
            w &= 5
    \end{work}
    In the story, that zero means:
    \ansline{after 5 washes the balance is \$0 --- the card is used up.}
    \ansline{A 6th wash would need more money.}
  \item Which inputs $w$ make sense here? \ans{whole washes, $w=0,1,2,3,4,5$} \; Is $B$ increasing or decreasing
        on them? \ans{decreasing}
  \item Suppose each wash cost \$10 instead of \$8. Would the card run out \emph{sooner} or
        \emph{later}? \ans{sooner ($w=4$)} \; On the graph, what changes --- where the line
        \emph{starts}, how \emph{steep} it is, or both? Say why.
        \ansline{Only the steepness --- the card still starts at \$40, so $(0,40)$ does not move;}
        \ansline{the slope becomes $-10$, a steeper drop, so the line reaches zero sooner.}
\end{enumerate}
\end{scenariobox}

\end{document}
